\documentclass[11pt,reqno]{amsart}
\usepackage[T1]{fontenc}
\usepackage{lmodern,amssymb,microtype}
\usepackage[colorlinks=true,linkcolor=black,citecolor=blue,urlcolor=black]{hyperref}
\hypersetup{pdftitle={A counterexample to the Erdos--Mahler conjecture},
  pdfsubject={Smooth convergents and prescribed irrationality exponents},
  pdfauthor={Diego Marques}}

\usepackage{enumitem}

\newtheorem*{theorem}{Theorem}
\newtheorem{lemma}{Lemma}[section]
\newtheorem{proposition}[lemma]{Proposition}
\theoremstyle{definition}
\newtheorem*{conjecture}{Conjecture}
\newcommand{\N}{\mathbb N}
\newcommand{\Q}{\mathbb Q}
\newcommand{\Z}{\mathbb Z}
\newcommand{\DD}{\mathcal D}
\newcommand{\PP}{\mathrm P}
\setlength{\emergencystretch}{1em}

\title{A counterexample to the Erd\H{o}s--Mahler conjecture}
\author[D. Marques]{Diego Marques}
\address{Departamento de Matem\'atica, Universidade de Bras\'ilia,
Bras\'ilia, DF, Brazil}
\email{diego@mat.unb.br}
\thanks{The author was supported by the Brazilian National Council for
Scientific and Technological Development (CNPq), Grant No.~304467/2023-5.}
\subjclass[2020]{Primary 11J70; Secondary 11J82, 11J86, 11J87}
\keywords{Erd\H{o}s--Mahler conjecture, continued fractions, smooth
convergents, irrationality exponent, Subspace Theorem, linear forms in
logarithms}

\begin{document}

\begin{abstract}
We disprove a 1939 conjecture of Erd\H{o}s and Mahler by constructing,
for every real $\tau\ge4$, a Cantor set of numbers of irrationality
exponent $\tau$, each having infinitely many regular continued-fraction
convergents of the form $2^a/3^b$.
\end{abstract}
\maketitle

\section{Introduction}\label{sec:intro}

For an irrational real number $\xi$, its \emph{irrationality exponent}
$\mu(\xi)$ is the supremum of the real numbers $\nu>0$ for which
$0<|\xi-u/v|<v^{-\nu}$ holds for infinitely many reduced fractions
$u/v$ with $v\ge1$. Thus $\xi$ is a Liouville number precisely when
$\mu(\xi)=\infty$. Write $p_n/q_n$ for the regular continued-fraction
convergents of $\xi$, and $\PP(m)$ for the largest prime factor of
$|m|$ when $m$ is a nonzero integer, with $\PP(\pm1)=1$.

In 1939, Erd\H{o}s and Mahler~\cite{EM} investigated
the prime factors of continued-fraction numerators
and denominators. They proved that if
$\PP(q_{n-1}q_nq_{n+1})$ is bounded for infinitely
many $n$, then $\xi$ is a Liouville number.
They conjectured that the same conclusion should
hold if $\PP(p_nq_n)$ is bounded for infinitely
many $n$. The question was subsequently studied by
Fraenkel~\cite{Fraenkel} and recalled by Erd\H{o}s
in his reminiscences of Mahler~\cite{ErdosReminiscences}. 
In 2004, Bugeaud~\cite[Problem~43, p.~215]{Bugeaud}
included the conjecture among the open problems
on continued fractions.

\begin{conjecture}[Erd\H{o}s--Mahler]
If $\PP(p_nq_n)$ is bounded for infinitely many $n$, then $\xi$ is
a Liouville number.
\end{conjecture}

We disprove the conjecture even when the selected numerators are
powers of $2$ and the corresponding denominators are powers of $3$.

\begin{theorem}
For every real $\tau\ge4$, there exists a Cantor set
$K_\tau\subset(1,2)$ such that every $\xi\in K_\tau$
has irrationality exponent $\tau$ and infinitely many regular
continued-fraction convergents of the form $2^a/3^b$, with $a,b$
positive integers.
\end{theorem}

By Ridout's theorem~\cite{Ridout}, bounded prime factors of the
denominators of infinitely many convergents already imply
transcendence. Fraenkel~\cite{Fraenkel62,Fraenkel} studied rational
approximation with prescribed prime factors and constructed Liouville
numbers with infinitely many convergents subject to such restrictions.
Fraenkel and Borosh~\cite{FraenkelBorosh} subsequently obtained
non-Liouville examples by allowing additional cofactors in the
numerators and denominators. When both cofactors are bounded, their
Hausdorff-dimension formula gives zero. 
Lelis and Marques~\cite[Theorem~1.1]{LelisMarques}
proved that $\xi$ is a Liouville number if there is an
increasing sequence $(n_j)$ such that
$\PP(p_{n_j}q_{n_j})$ is bounded and
$n_{j+1}-n_j=o(\log q_{n_j})$ as $j\to\infty$.


More recently, Bugeaud and Nguyen~\cite[Question~4.1, p.~32]{BN}
revisited the conjecture and constructed further examples with
restricted prime factors, leaving open whether the numbers they
constructed are Liouville. This illustrates the main difficulty:
prescribing infinitely many convergents does not exclude
exceptionally good rational approximations at other denominators.

Our construction selects successive fractions $2^a/3^b$
in short intervals to the right of the preceding choice.
Their prescribed displacement gives the lower bound
$\mu(\xi)\ge\tau$; the difficulty is to exclude better
rational approximations at other denominators. For
$\nu>4$, the neighborhoods of radius $v^{-\nu}$ around
reduced fractions in $[1,2]$ with denominator $v>D$
have a union of measure $O(D^{2-\nu})$. Proposition~\ref{prop:count} gives the same upper bound,
with an absolute implied constant, for the upper limiting
proportion of prescribed fractions admitting such competing
approximations. The estimate holds despite their rigid
arithmetic structure and is independent of the intervals
used in the construction.

The proof begins with an elementary counting estimate for
irrational rotations, which handles small denominators.
For the remaining approximations, the $p$-adic Subspace
Theorem improves the elementary bound $\log v=O(N)$ to
$\log v=o(N)$ on blocks of $N$ indices. The main
combinatorial observation is that, for a fixed sum of
two indices, the product of the corresponding fractions
takes at most two values. In denominator ranges
$(V,V^\theta]$ with $1<\theta<\nu/4$, these identities
persist for sufficiently close rational approximants,
whose product set consequently has cardinality $O(N)$.
The Pl\"unnecke--Ruzsa inequality then bounds the number
of classes these approximants occupy modulo $q$th powers,
for a suitable integer $q$. A lower bound for linear forms
in logarithms, together with $3$-adic valuations, confines
the corresponding indices in each class to a short
interval. The bound $\log v=o(N)$ makes the resulting
estimates summable over all denominator ranges.


\section{The counting proposition and the main theorem}
\label{sec:main}

We first state the counting proposition and deduce the main theorem from it. The proposition is proved in Section~\ref{sec:count}.

All logarithms are natural. We write $\N=\{1,2,\ldots\}$,
take rational fractions with positive denominators, and
use $|I|$ for the length of an interval or the cardinality
of a finite set. Put
\[
 \alpha=\frac{\log 3}{\log 2},\qquad
 a(b)=\lceil b\alpha\rceil,\qquad
 R(b)=\frac{2^{a(b)}}{3^b}\qquad(b\in\N).
\]
Since $\alpha$ is irrational, $0<a(b)-b\alpha<1$.
Thus $R(b)$ is the unique fraction in $(1,2)$ with
denominator $3^b$ and numerator a power of~$2$.
It is in lowest terms. Let $\DD$ be the set of
denominators of the regular continued-fraction
convergents of $\alpha$.

At each stage of the construction, we shall select fractions
$R(b)$ lying in a prescribed short interval. We must also
exclude those admitting exceptionally close competing rational
approximations. The following proposition controls these
exceptions over the entire block $[N,2N)$, independently
of the target interval. Crucially, its constant is independent
of the approximation exponent, which will approach $4$
along the construction.

Fix $4<\nu\le5$. A reduced fraction $u/v\in[1,2]$ is
\emph{exceptional for $b$ at exponent $\nu$} if
$u/v\ne R(b)$ and
\begin{equation}\label{eq:exception}
 \left|R(b)-\frac uv\right|<\frac1{16v^\nu}.
\end{equation}
For $D\ge2$, let $E_N(D,\nu)$ count the integers
$b\in[N,2N)$ for which there is an exceptional
fraction with $v>D$.

\begin{proposition}[Exceptional approximations]
\label{prop:count}
There is an absolute constant $C$ such that, for every
$4<\nu\le5$ and every $D\ge2$,
\[
 \limsup_{\substack{N\to\infty\\N\in\DD}}
 \frac{E_N(D,\nu)}{N}\le CD^{2-\nu}.
\]
\end{proposition}

We now fix a real $\tau\ge4$ for the construction.
For $t\in\mathbb N$ put $Q=3^t$ and define
\[
 \begin{aligned}
 J_t&=(R(t)+Q^{-\tau},R(t)+2Q^{-\tau}),\\
 K(t)&=[R(t)+Q^{-\tau},R(t)+3Q^{-\tau}].
 \end{aligned}
\]
The next fraction will be chosen in $J_t$, while $K(t)$
will contain the corresponding limit. Both intervals lie
in $(1,2)$ for sufficiently large $t$, since
$R(t)\le2-Q^{-1}$.

Put
\[
 D_t=Q^{\tau-1}/6,\qquad \nu_t=4+1/t.
\]
An index $b$ is \emph{admissible at $t$} if $R(b)\in J_t$
and $b$ has no exceptional fraction at exponent $\nu_t$
with denominator $v>D_t$. The threshold $D_t$ separates the denominators handled
by rational separation from those controlled by admissibility. 
For these parameters, Proposition~\ref{prop:count} bounds
the upper limiting proportion of exceptional indices by
$O(Q^{-2(\tau-1)})$, whereas the proportion of indices $b$
with $R(b)\in J_t$ tends to a quantity of order $Q^{-\tau}$.
Thus sufficiently many admissible indices remain.

\begin{lemma}\label{lem:choice}
For every sufficiently large $t$, all sufficiently large $N\in\DD$
admit two distinct indices $b,c\in[N,2N)$ admissible at $t$.
\end{lemma}

\begin{proof}
Fix $t$ and put
\[
 \Phi(y)=1-\frac{\log y}{\log2}
 \quad(1\le y\le2),\qquad
 \rho_t=|\Phi(J_t)|.
\]
Write $\{x\}=x-\lfloor x\rfloor$ for fractional part. Since
$\Phi(R(b))=\{b\alpha\}$ and $|\Phi'(y)|=1/(y\log2)$, we have
\[
 \frac{Q^{-\tau}}{2\log2}\le\rho_t\le\frac{Q^{-\tau}}{\log2}.
\]
By the grid estimate \eqref{eq:grid}, proved at the beginning
of Section~\ref{sec:count}, there are $N\rho_t+O(1)$ indices
with $R(b)\in J_t$. Among these we discard at most $E_N(D_t,\nu_t)$,
since this last count includes exceptional indices in the whole of
$[N,2N)$. Writing $A_t(N)$ for the number of admissible indices
that remain, we obtain from Proposition~\ref{prop:count}
\[
 \liminf_{\substack{N\to\infty\\N\in\DD}}\frac{A_t(N)}{N}
 \ge \rho_t-C D_t^{2-\nu_t}\ge\rho_t-36C Q^{-2(\tau-1)}.
\]
Here $D_t^{2-\nu_t}\le D_t^{-2}$, and the last error divided
by $\rho_t$ is $O(Q^{2-\tau})$. Hence the right-hand side exceeds
$\rho_t/2$ for every sufficiently large $t$. With such a $t$ fixed,
$A_t(N)$ therefore tends to infinity.
\end{proof}

\begin{proof}[Proof of the theorem]
For any $t$ large enough for Lemma~\ref{lem:choice}, take
$N\in\DD$ larger than $Q=3^t$ and sufficiently large to provide
two admissible indices $b<c$ in $[N,2N)$. We verify that their
intervals lie strictly inside $K(t)$ and are disjoint.

For $d=b$ or $c$, the left endpoint of $K(d)$ exceeds $R(t)+Q^{-\tau}$.
At the other endpoint, $R(d)\in J_t$ and $3^d\ge2Q$ give
\[
 R(d)+3\cdot3^{-\tau d}
 <R(t)+(2+3\cdot2^{-\tau})Q^{-\tau}<R(t)+3Q^{-\tau}.
\]
To see disjointness, use the fact that the two indices were taken
from the same interval $[N,2N)$, so that $c<2b$. The distinct
fractions $R(b)$ and $R(c)$ can both be written with denominator
$3^c$, and hence differ by at least $3^{-c}>3^{-2b}$. Each of the
two intervals lies within $3\cdot3^{-\tau b}<3^{-2b}$ to the right of
its corresponding fraction.

Starting at a sufficiently large $t_0$, repeat the selection in
each interval, labeling the two choices $0$ and $1$ in order of
their indices. An infinite binary sequence determines a nested
chain of intervals $K(t_j)$. Write $Q_j=3^{t_j}$ and $r_j=R(t_j)$
along this chain. Since $Q_{j+1}\ge2Q_j$, every interval at level
$j$ has length at most $2Q_0^{-\tau}2^{-\tau j}$. The chain therefore
has a unique common point $\xi$. Sequences agreeing through level $j$
give points in the same interval at that level, and sequences making
different choices give distinct points by disjointness. Thus the map
from $\{0,1\}^{\N}$ to its image is continuous and injective. By
compactness it is a homeomorphism onto a Cantor set $K_\tau\subset(1,2)$.

Fix any $\xi\in K_\tau$ and consider its chain. Since the next interval
lies in the interior of the preceding one, we have
\begin{equation}\label{eq:displacement}
 Q_j^{-\tau}<\xi-r_j<3Q_j^{-\tau}\qquad(j\ge0).
\end{equation}
If $\xi=u/v$ were rational, its distance from $r_j$ would be at
least $1/(vQ_j)$, in contradiction with \eqref{eq:displacement} for
large $j$. Moreover, $3Q_j^{-\tau}<1/(2Q_j^2)$, so Legendre's
criterion \cite[Theorem~1.8, p.~10]{Bugeaud} makes each $r_j$ a
regular convergent of $\xi$. This gives infinitely many convergents
$2^{a(t_j)}/3^{t_j}$ since the denominators strictly increase.
For every $s<\tau$, their errors are less
than $Q_j^{-s}$ for all sufficiently large $j$, which gives
$\mu(\xi)\ge\tau$.

For the upper bound, let $u/v\in[1,2]$ be reduced with $v\ge Q_0$.
Choose the unique $j$ for which $Q:=Q_j\le v<Q_{j+1}=:S$.
For $v\le D_{t_j}$ we compare $u/v$ with $r_j$. If $u/v=r_j$,
then $v=Q$ and \eqref{eq:displacement} gives $|\xi-u/v|>v^{-\tau}$.
Otherwise rational separation gives
\[
 |\xi-u/v|\ge\frac1{vQ}-\frac3{Q^\tau}
 \ge\frac1{2vQ}\ge\frac1{2v^2},
\]
where we used $6v\le Q^{\tau-1}$ and $Q\le v$.

For $D_{t_j}<v<S$ we compare instead with $r_{j+1}$. The condition
$v<S$ ensures that $u/v\ne r_{j+1}$, and gives
\[
 |r_{j+1}-u/v|\ge\frac1{Sv}>S^{-2}.
\]
Since $v<S$ and $S\ge3$, we have $6v<S^3$, and hence
\[
 |\xi-r_{j+1}|<\frac3{S^\tau}\le\frac3{S^4}
 \le\frac1{2Sv}
 \le\frac12|r_{j+1}-u/v|.
\]
As $v>D_{t_j}$, the admissibility of $t_{j+1}$ at
$t_j$ gives
\[
 |\xi-u/v|\ge\tfrac12|r_{j+1}-u/v|
 \ge \frac1{32v^{\nu_{t_j}}}.
\]

For $Q_j\le v<Q_{j+1}$ and $u/v\ne r_j$, the two ranges give
$|\xi-u/v|\ge(32v^{\nu_{t_j}})^{-1}$. Given $\varepsilon>0$, we have
$\nu_{t_j}\le4+\varepsilon/2$ for all sufficiently large $j$.
For all sufficiently large $v$ we also have $v^{\varepsilon/2}\ge32$,
so $|\xi-u/v|\ge v^{-4-\varepsilon}$ for all such fractions.
For the prescribed convergents themselves, \eqref{eq:displacement}
gives $|\xi-r_j|>Q_j^{-\tau}$. Since $\tau\ge4$ and rationals
outside $[1,2]$ are bounded away from $\xi$, it follows that
$\mu(\xi)\le\tau$. Thus $\mu(\xi)=\tau$ for every $\xi\in K_\tau$.
\end{proof}

\section{Proof of Proposition~\ref{prop:count}}\label{sec:count}

We begin with the elementary rotation estimate used in
Lemma~\ref{lem:choice}. For every arc
$I\subset\mathbb R/\mathbb Z$ and $N\in\DD$,
\begin{equation}\label{eq:grid}
 \#\{N\le b<2N:\{b\alpha\}\in I\}
 =N|I|+O(1),
\end{equation}
with an absolute implied constant. To prove this, let
$p/N$ be the convergent of $\alpha$ with denominator $N$.
Since $(p,N)=1$, the points
\[
 N\alpha+\frac{ip}{N}\pmod 1,
 \qquad 0\le i<N,
\]
form a translate of the uniform grid of spacing $1/N,
$ which contains $N|I|+O(1)$ points in $I$.
By~\cite[Corollary~1.4, p.~10]{Bugeaud},
we have $|\alpha-p/N|<N^{-2}$. Thus replacing these
points by $(N+i)\alpha$ moves each by less than $1/N$.
Only grid points within this distance of an endpoint
of $I$ can change membership, and there are $O(1)$ such points.

Fix $4<\nu\le5$ and let $N$ tend to infinity through $\DD$.
We divide the exceptional approximations according as their
denominators are at most $N^{1/3}$ or exceed it. In the first range
we can sum the error term in \eqref{eq:grid} over all the fractions
in question. In the second, the interval belonging to a single
fraction has length smaller than the grid spacing $1/N$. Thus each
fraction occurs for only a bounded number of indices. The problem
in this range is to bound the number of distinct fractions that occur.

\subsection{The range of exceptional denominators}

For $V\ge2$, there are $O(V^2)$ reduced fractions in $[1,2]$ with
denominator in $(V,2V]$. For such a fraction $u/v$, the image under
$\Phi$ of its $(16v^\nu)^{-1}$-neighborhood in $[1,2]$ is an
interval of length $O(V^{-\nu})$. The grid estimate \eqref{eq:grid}
therefore gives $O(NV^{-\nu}+1)$ possible indices for this fraction.
Summing first over the fractions and then over $V=2^iD$, $i\ge0$,
with $V<N^{1/3}$, we bound the number of indices admitting an
exceptional denominator in $(D,N^{1/3}]$ by an absolute constant times
\begin{equation}\label{eq:small}
 \sum_{V}\bigl(NV^{2-\nu}+V^2\bigr)
 \ll ND^{2-\nu}+N^{2/3}.
\end{equation}
The decreasing sum is bounded in terms of its first term, since
$2^{2-\nu}\le1/4$. The increasing sum is $O(N^{2/3})$, so the accumulated grid error
is $o(N)$.

For larger denominators, rational separation alone gives
$16v^{\nu-1}<3^b$ by \eqref{eq:exception}. In particular it gives
only $\log v=O(N)$ on $[N,2N)$. To sum the estimates for separate
denominator ranges we shall need $\log v=o(N)$ uniformly over the
exceptional fractions. The following consequence of the $p$-adic
Subspace Theorem supplies this improvement.

\begin{lemma}\label{lem:height}
For every $\varepsilon>0$, all sufficiently large $b$ and
all reduced $u/v\in[1,2]$ distinct from $R(b)$ satisfy
\[
 |R(b)-u/v|\ge 3^{-\varepsilon b}v^{-2}.
\]
\end{lemma}

\begin{proof}
Write $P=2^{a(b)}$ and $S=3^b$. If $P/S$ is close to $u/v$, the
integers $Pv$ and $Su$ are close relative to their size. We apply
the Subspace Theorem to the primitive pair with this ratio. Put
\[
 d=\gcd(Pv,Su),\qquad U=Pv/d,\qquad W=Su/d,\qquad e=|Pv-Su|.
\]
We have $e\ge1$. Every prime divisor of $d$ belongs to $\{2,3\}$,
since any other common prime divisor would divide both $u$ and $v$.
For a positive integer $n$, put $[n]_6=2^{v_2(n)}3^{v_3(n)}$, where
$v_p$ is the usual $p$-adic valuation. With $|p|_p=p^{-1}$ we have
\[
 |UW|_2|UW|_3=[UW]_6^{-1}=\frac{d^2}{PS[uv]_6}.
\]
At the real place we take the forms $X-Y,Y$, so the factor in the
product there is $|(U-W)W|=eSu/d^2$. Taking $X,Y$ at each of
$2$ and $3$ gives
\[
 \Pi:=|(U-W)W|\,|UW|_2|UW|_3
 =\frac{eu}{P[uv]_6}
 \le 2v^2|P/S-u/v|.
\]
For the last inequality, use $e=Sv|P/S-u/v|$, $u\le2v$, and $S<P$.

Suppose now that the asserted inequality fails along a sequence
$b\to\infty$. On this sequence
\[
 (Sv)^{-1}\le |P/S-u/v|<S^{-\varepsilon}v^{-2},
\]
so $v<S$. It follows that $H=\max(U,W)\le2Sv<2S^2$ and hence
\[
 \Pi<2S^{-\varepsilon}\le2^{1+\varepsilon/2}H^{-\varepsilon/2}.
\]
The forms at each of the three places have determinant one.
The $p$-adic Subspace Theorem
\cite[p.~246, (1.3)--(1.5)]{Schlickewei} therefore implies that
all such pairs of sufficiently large height lie on finitely many
rational lines through the origin.
Indeed, the right-hand side is smaller than $(U^2+W^2)^{-\delta/2}$
for all sufficiently large $H$, with
$\delta=\min(\varepsilon/4,1/2)>0$.

Each rational line contains at most one positive primitive
pair, and there are only finitely many pairs of bounded
height. Thus $U/W$ belongs to a finite set not containing
$1$. On the other hand,
\[
 0<|U/W-1|
 =\frac{|P/S-u/v|}{u/v}
 <S^{-\varepsilon}v^{-2}
 \le 3^{-\varepsilon b}\longrightarrow0
\]
as $b\to\infty$, a contradiction.
\end{proof}

Let $W_N$ be the largest value of $\log v$ among exceptional
fractions for $N\le b<2N$, taking $W_N=0$ when there are none.
The preceding rational separation bound ensures that the maximum
is finite. For each $\varepsilon>0$, Lemma~\ref{lem:height} and
\eqref{eq:exception} give $16v^{\nu-2}<3^{\varepsilon b}$ as
soon as $N$ is sufficiently large. Since $b<2N$,
\[
 \limsup_{N\to\infty,\ N\in\DD}\frac{W_N}{N}
 \le\frac{2\varepsilon\log3}{\nu-2}.
\]
Letting $\varepsilon$ decrease to zero, we conclude that
\begin{equation}\label{eq:height}
 W_N=o(N).
\end{equation}

\subsection{Products in a fixed denominator range}

Choose $\theta$ with $1<\theta<\nu/4$. The condition
$4\theta<\nu$ ensures that the error in a product
of two exceptional approximants is smaller than
the separation between distinct rational products.
For $w>0$ put $H=e^w$ and $V=e^{w/\theta}$,
so that $H=V^\theta$. Write
$\mathcal A(w)$ for the set of indices $b\in[N,2N)$ admitting an
exceptional fraction with $v>N^{1/3}$ and $V<v\le H$. We shall
first estimate $|\mathcal A(w)|$ for a single $w$, and only at
the end choose the ranges over which to sum.

Assume $\mathcal A(w)$ is nonempty. For each of its indices choose
one such fraction $s_b$, and let $X$ be the set of distinct fractions
so obtained. Write $M=|X|$. If the same fraction of denominator $v$
is chosen several times, \eqref{eq:grid} bounds the number of those
indices by $O(1+Nv^{-\nu})=O(1)$, since $v>N^{1/3}$ and $\nu>4$.
Consequently
\begin{equation}\label{eq:multiplicity}
 |\mathcal A(w)|\ll M .
\end{equation}
We may henceforth retain one index for each element of $X$.

For these indices, compare the products $s_bs_c$ with the fractions
\[
 R(b)R(c)=2^{a(b)+a(c)}3^{-(b+c)} .
\]
For fixed $\ell=b+c$, the integer $a(b)+a(c)$ lies in
$(\alpha\ell,\alpha\ell+2)$ and thus takes at most two values.
There are fewer than $2N$ choices for $\ell$. It follows that the
products $R(b)R(c)$ take at most $4N$ values.

The choice of $\theta$ allows us to transfer this bound to
$XX=\{ss':s,s'\in X\}$. Indeed, \eqref{eq:exception} and the
fact that all fractions lie in $[1,2]$ give
\[
 |s_bs_c-R(b)R(c)|
 \le 2|s_b-R(b)|+2|s_c-R(c)|<\tfrac14 V^{-\nu}.
\]
If $R(b)R(c)=R(b')R(c')$, the two corresponding products of elements
of $X$ thus differ by less than $\tfrac12 V^{-\nu}$. Their reduced
denominators are at most $H^2$, so distinct products would differ by
at least $H^{-4}$. Since $H^4=V^{4\theta}<V^\nu$, the products
are equal. We have proved
\begin{equation}\label{eq:products}
 |XX|\le4N .
\end{equation}

When $M$ is a substantial fraction of $N$, \eqref{eq:products}
bounds the ratio $|XX|/M$. We use this ratio to bound the number of
classes occupied by $X$ modulo rational $q$th powers. Two positive
rationals belong to the same class when their quotient is a $q$th
power in $\Q$. For finite subsets $A,B$ of an abelian group, write
$A\pm B=\{a\pm b:a\in A,\ b\in B\}$ and
$m\cdot A=\{ma:a\in A\}$. In the next lemma
the group law is written additively.

\begin{lemma}\label{lem:classes}
Let $A$ be a finite nonempty subset of a torsion-free abelian group
$G$, and put $L=|A+A|/|A|$. For $q=2^k$ with $k\ge1$, the image
of $A$ in $G/qG$ has at most $L^{3k}$ elements.
\end{lemma}

\begin{proof}
Choose one representative $a\in A$ of each class in the image of
$A$. The sets $a-q\cdot A$ lie in distinct cosets of $qG$ and
each contains $|A|$ elements, since dilation by $q$ is injective.
Their union is contained in $A-q\cdot A$. It therefore suffices
to bound $|A-q\cdot A|/|A|$.

For a single doubling, the Pl\"unnecke--Ruzsa inequality
\cite[Theorem~1.2]{Petridis}, applied to $A$ and itself, gives
\[
 |A-2\cdot A|\le|A-A-A|\le L^3|A|.
\]
For the iteration we use Ruzsa's triangle inequality
$|U-W|\,|V|\le|U-V|\,|V-W|$ for finite nonempty sets
\cite[proof of Theorem~1.2, (3)]{Petridis}. Take
$U=A$, $V=2^{i-1}\cdot A$, $W=2^i\cdot A$. Since
$V-W=2^{i-1}\cdot(A-2\cdot A)$, injectivity of dilation gives
\[
 \frac{|A-2^i\cdot A|}{|A|}
 \le\frac{|A-2^{i-1}\cdot A|}{|A|}
     \frac{|A-2\cdot A|}{|A|}\qquad(i\ge2).
\]
Iterating the single-doubling bound proves the lemma.
\end{proof}

\subsection{The count at large denominators}

By Lemma~\ref{lem:classes} in $\Q_{>0}^{\times}$ and
\eqref{eq:products}, the elements of $X$ occupy at most
$(4N/M)^{3k}$ classes modulo rational $q$th powers when $q=2^k$.
To bound the size of each class we choose $k$ in terms of
$x=\log(N/w)$. Here $N$ is the length of the interval of indices
and $w$ is the upper bound on $\log v$ in the denominator range.

\begin{lemma}\label{lem:block}
There are absolute constants $c,C,x_0>0$ such that,
for $4<\nu\le5$, $1<\theta<\nu/4$, $N\in\DD$,
and $w\ge1$ with $x=\log(N/w)\ge x_0$,
\[
 \frac{|\mathcal A(w)|}{N}
 \le C\exp\left(-\frac{cx}{\log(2+x)}\right).
\]
\end{lemma}

\begin{proof}
We may suppose that $\mathcal A(w)$ is nonempty and $x\ge2$.
Let $q=2^k$ be the least power of two with $q\ge x^2$, so that
$x^2\le q<2x^2$. We shall show that within a class modulo $q$th
powers the ratio of two exceptional fractions agrees with the ratio
of the corresponding $R(b)$. The height of the first ratio will
then bound the difference of the two indices.

Take distinct $s_b,s_c$ in the same class and write $s_c/s_b=g^q$
with $g\in\Q_{>0}$, $g\ne1$. For a positive rational $r/s$ in
lowest terms, write $h(r/s)=\log\max(r,s)$. Each of $s_b,s_c$
has denominator at most $H$ and numerator at most $2H$. Thus the
numerator and denominator of their quotient in lowest terms are
both at most $2H^2$. Taking the $q$th root gives
\begin{equation}\label{eq:ratioheight}
 qh(g)=h(s_c/s_b)\le 2w+\log2\le 3w .
\end{equation}
This gain of a factor $q$ in the height is the reason for working
within one power class.

Consider the difference of the logarithms of the two ratios,
\[
 \Lambda=(a(c)-a(b))\log2-(c-b)\log3-q\log g.
\]
The logarithm has derivative at most one on $[1,2]$. Comparing
$\log s_b$ with $\log R(b)$, and likewise at $c$, we obtain
from \eqref{eq:exception}
\begin{equation}\label{eq:logupper}
 |\Lambda|<\tfrac18 V^{-\nu}=\tfrac18\exp(-\nu w/\theta).
\end{equation}

Suppose $\Lambda\ne0$. We use the weighted coefficient parameter
in Matveev's lower bound
\cite[Corollary~2.3, (2.6), with (1.3)]{Matveev}.
For positive rational bases $\gamma_i\ne1$ and integer coefficients
$m_i$ $(1\le i\le3)$, take
\[
 A_i\ge\max\{h(\gamma_i),|\log\gamma_i|,0.16\},
 \qquad B=\max\left\{1,\max_i\frac{|m_i|A_i}{A_3}\right\}.
\]
For a nonzero form in their real logarithms the bound over $\Q$ is
\[
 \log\left|\sum_{i=1}^3m_i\log\gamma_i\right|
 >-C_0 A_1A_2A_3(1+\log B)
\]
with an absolute $C_0$. By \eqref{eq:ratioheight},
the existence of distinct
$s_b,s_c$ in the same class forces
$3w/q\ge h(g)\ge\log2$. Thus we may take
\[
 A_1=\log2,\qquad A_2=\log3,\qquad A_3=3w/q,
\]
since $|\log g|\le h(g)$. The cited bound requires
only that the linear form be nonzero; no multiplicative
independence of $2,3$, and $g$ is assumed.

Since $b,c\in[N,2N)$, we have
$|c-b|<N$ and $|a(c)-a(b)|\le\alpha N+1$.
Thus the coefficients of $\log2$ and $\log3$
are $O(N)$, whereas that of $\log g$ is $q$.
With $A_3=3w/q$, the weighted definition of $B$
gives $B\ll qN/w$ and hence
$1+\log B\ll1+\log q+x$. Substitution in Matveev's bound gives
\[
 -\log|\Lambda|\ll\frac{w}{q}(1+\log q+x)\ll\frac wx .
\]
This is where the choice $q\asymp x^2$ enters. On the other hand,
\eqref{eq:logupper} gives $-\log|\Lambda|>(\nu/\theta)w>4w$.
The two bounds are incompatible for $x$ larger than an absolute
constant. It follows that $\Lambda=0$, as claimed.

The vanishing of $\Lambda$ turns the rational-power
relation into an exact identity between the prescribed
fractions:
\[
 \frac{s_c}{s_b}=2^{a(c)-a(b)}3^{b-c}.
\]
Taking the valuation at $3$ and using \eqref{eq:ratioheight} gives
$|b-c|\log3\le h(s_c/s_b)\le2w+\log2$. Once $s_b$ is fixed, all
the representative indices in its class consequently lie in an
interval of length $O(w)$. They are distinct integers, so the class
has $O(w)$ elements, with an absolute constant.

Multiplying by the bound for the number of classes gives
\[
 M\le C_1 w(4N/M)^{3k}.
\]
To read this as a bound for the proportion of exceptional indices,
put $y=M/(4N)$. Since $w/N=e^{-x}$, the preceding inequality gives
$y^{3k+1}\le(C_1/4)e^{-x}$. Although $q$ is of order $x^2$,
only $k=\log_2 q\ll\log(2+x)$ enters this exponent. Taking the
$(3k+1)$st root and applying \eqref{eq:multiplicity} proves the
assertion, after increasing $x_0$ if necessary.
\end{proof}

\begin{proof}[Completion of the proof of Proposition~\ref{prop:count}]
It remains to sum the estimates over the denominator
ranges. Their number need not be bounded; what matters
is that \eqref{eq:height} forces all nonempty ranges
into the tail of a convergent series. Put
\[
 w_j=N\theta^{-j},\qquad x_j=\log(N/w_j)=j\log\theta
 \qquad(j=0,1,2,\ldots).
\]
The intervals $(w_j/\theta,w_j]$ partition $(0,N]$. By
\eqref{eq:height}, they contain the logarithms of all exceptional
denominators exceeding $N^{1/3}$ once $N$ is sufficiently large.
For any $j$ with $\mathcal A(w_j)\ne\varnothing$, an exceptional
denominator in that range satisfies
\[
 \tfrac13\log N<\log v\le w_j<\theta\log v\le\theta W_N .
\]
Thus $w_j\ge1$ for large $N$. Moreover, if $W_N>0$, then
$x_j>\log(N/(\theta W_N))$, which tends to infinity by
\eqref{eq:height}. This places every nonempty range in the part of
the estimate where Lemma~\ref{lem:block} applies.

For fixed $\theta>1$, the series
\[
 \sum_{j\ge0}\exp\!\left(-\frac{c x_j}{\log(2+x_j)}\right)
\]
converges, since for large $j$ its terms are at most
$\exp(-c_\theta\sqrt j)$ for some $c_\theta>0$. Given any tail
of this series, the nonempty ranges all belong to that tail for
sufficiently large $N$. Lemma~\ref{lem:block} therefore gives
$\sum_j|\mathcal A(w_j)|=o(N)$. An index may occur in several
ranges, but this sum bounds the number admitting any exceptional
denominator above $N^{1/3}$.

Combining with \eqref{eq:small} yields
\[
 E_N(D,\nu)\le C\bigl(ND^{2-\nu}+N^{2/3}\bigr)+o_\nu(N),
\]
where $C$ is absolute and $\theta$ has been fixed in terms of $\nu$.
Only the last term uses convergence for this particular $\theta$.
Dividing by $N$ and taking the upper limit along $\DD$ proves the
proposition.
\end{proof}

\section*{AI use declaration}

The author used ChatGPT and Codex in developing the counting
argument of Section~3, including the reduction to twofold
products and the use of rational-power classes, and in
formulating the diagonal construction in Section~2.
These tools were also used to review the exposition and
check proofs and references. The author takes full
responsibility for the mathematical arguments and citations.

\begin{thebibliography}{99}

\bibitem{Bugeaud}
Y. Bugeaud, \emph{Approximation by Algebraic Numbers}, Cambridge
Tracts in Mathematics \textbf{160}, Cambridge University Press,
Cambridge, 2004.

\bibitem{BN}
Y. Bugeaud and K. D. Nguyen,
\href{https://doi.org/10.2140/pjm.2023.326.17}{Some arithmetical
properties of convergents to algebraic numbers},
\emph{Pacific J. Math.} \textbf{326} (2023), no.~1, 17--36.

\bibitem{EM}
P. Erd\H{o}s and K. Mahler,
\href{https://doi.org/10.1112/jlms/s1-14.1.12}{Some arithmetical
properties of the convergents of a continued fraction},
\emph{J. London Math. Soc.} \textbf{14} (1939), no.~1, 12--18.

\bibitem{ErdosReminiscences}
P. Erd\H{o}s,
\href{https://old.renyi.hu/~p_erdos/1989-34.pdf}
{Some personal and mathematical reminiscences
of Kurt Mahler},
\emph{Austral. Math. Soc. Gaz.}
\textbf{16} (1989), no.~1, 1--2.

\bibitem{Fraenkel62}
A. S. Fraenkel,
\href{https://doi.org/10.1090/S0002-9947-1962-0144852-2}{On a
theorem of D. Ridout in the theory of Diophantine approximations},
\emph{Trans. Amer. Math. Soc.} \textbf{105} (1962), 84--101.

\bibitem{Fraenkel}
\bysame,
\href{https://doi.org/10.1112/jlms/s1-39.1.405}{Transcendental
numbers and a conjecture of Erd\H{o}s and Mahler},
\emph{J. London Math. Soc.} \textbf{39} (1964), 405--416.

\bibitem{FraenkelBorosh}
A. S. Fraenkel and I. Borosh,
\href{https://doi.org/10.1112/plms/s3-15.1.458}{Fractional dimension
of a set of transcendental numbers}, \emph{Proc. London Math. Soc.}
(3) \textbf{15} (1965), 458--470. Corrigendum, \emph{ibid.}
(3) \textbf{16} (1966), 192.

\bibitem{LelisMarques}
J. Lelis and D. Marques,
\href{https://doi.org/10.1017/S0004972716000770}
{On a problem of Erd\H{o}s and Mahler concerning
continued fractions},
\emph{Bull. Aust. Math. Soc.}
\textbf{95} (2017), no.~2, 183--186.

\bibitem{Matveev}
E. M. Matveev,
\href{https://doi.org/10.1070/IM2000v064n06ABEH000314}{An explicit
lower bound for a homogeneous rational linear form in the logarithms
of algebraic numbers. II}, \emph{Izv. Math.} \textbf{64} (2000),
no.~6, 1217--1269.

\bibitem{Petridis}
G. Petridis,
\href{https://doi.org/10.1007/s00493-012-2818-5}{New proofs of
Pl\"unnecke-type estimates for product sets in groups},
\emph{Combinatorica} \textbf{32} (2012), no.~6, 721--733,
\href{https://arxiv.org/abs/1101.3507v3}{arXiv:1101.3507v3}.

\bibitem{Ridout}
D. Ridout,
\href{https://doi.org/10.1112/S0025579300001182}{Rational
approximations to algebraic numbers},
\emph{Mathematika} \textbf{4} (1957), 125--131.

\bibitem{Schlickewei}
H. P. Schlickewei,
\href{https://www.numdam.org/item/CM_1992__82_3_245_0/}{The
quantitative Subspace Theorem for number fields},
\emph{Compositio Math.} \textbf{82} (1992), no.~3, 245--273.

\end{thebibliography}
\end{document}
